% Research lemma; not a standalone article and not a kernel-checked proof.
\section{An unbounded-time obstruction for a read-only gap detector}
\label{sec:read-only-detector}

The following obstruction uses no running-time bound.  It applies only when
the detector writes no tape symbols and counts every internal control state.
An exit is terminal for purposes of the classifier: if a caller performs
further observations that can affect the classification, that caller's states
must be included in the state count.

\paragraph{Zero-excursion lemma.}
Let a deterministic read-only machine have $n\ge1$ internal states and unit
left/right moves.  Consider a maximal consecutive sequence of its internal
observations of symbol zero.  Suppose it begins at position $s$, and all tape
positions strictly below a first right nonzero barrier $b>s$ are zero.
If the excursion is finite, it observes no cell below $s-(n-1)$.
The excursion may end by reaching the right barrier or by leaving the internal
state set.  The position reached by a terminal exit move, which is not then
read by the detector, is not an observation in this statement.

\paragraph{Proof.}
On symbol zero, the internal state determines the next state and the signed
move, independently of position.  If the observed internal states do not
repeat, there are at most $n$ observations.  Their positions are reached after
at most $n-1$ moves from $s$, so the claimed bound holds.

Otherwise take the first repeated state: $q_j=q_i$, where $0\le i<j\le n$.
Every transition before this recurrence remains internal.  Let
$\Delta=p_j-p_i$ be the displacement of this deterministic zero-state cycle.
No position observed during the first completed cycle reaches $b$.
If $\Delta\le0$, each successive copy of the cycle is a translation of the
first by $k\Delta\le0$.  Such a copy cannot newly reach the right barrier.
It also cannot exit the state set, since all the repeated zero transitions are
internal.  Thus the excursion would be infinite, contrary to hypothesis.

Hence $\Delta>0$.  All later cycles lie strictly to the right of their
predecessors, so the smallest position of the entire excursion occurs in the
transient or first cycle.  Moreover $p_j=p_i+\Delta>p_i$, so that minimum
occurs at one of $p_0,\ldots,p_{j-1}$, reached in at most $n-1$ moves from $s$.
Its position is at least $s-(n-1)$, as claimed.  If a later positive cycle hits
the barrier before it completes, this can only shorten the excursion and
cannot decrease its minimum.

\paragraph{Application to the gap-7 versus long classification.}
Let the long input have a fixed one at $+2$ and zeros at all positions
$\le1$, and let the head begin at $0$.  Suppose the classifier halts on this
input.  Its first zero excursion starts at $s=0$, so with seven states its
observed minimum is at least $-6$.  Any later excursion entering this left
zero region must cross the fixed one at $+2$ and begin at $s=1$, so its minimum
is at least $-5$.  Computation to the right of the known one, however long or
complicated, does not alter either bound: the tape is read-only and every
return carries only an internal state.

Compare a short input that is identical to the long input at all positions
greater than $-7$, but has its first differing one at $-7$.  To distinguish
the two inputs, a deterministic classifier must observe a cell at which they
differ.  With unit moves, reaching any smaller differing position first
requires observing $-7$.  The terminating long computation never does so.
The two computations consequently have the same state, position, observation,
and transition at every step up to the common exit, and give the same answer.
They cannot be classified differently.

Therefore seven internal read-only states are insufficient, even with
unbounded running time, arbitrary right-side computation, and repeated visits
to the known one.  The conclusion is stronger than a bounded search result for
a read-only seven-state template.  It does not exclude a detector that writes
temporary tape marks, shares additional states with an interpreter, or returns
a tape/head condition whose interpretation uses uncounted caller states.
